% Beyond 1.58 Bits: Adaptive-Precision Ternary Transformers
% with Fully-Ternary Mixture-of-Experts Routing
% Author: CrashOverrideX, Quillan AI Research Team
% Overleaf-ready: paste into a blank Overleaf project (pdfLaTeX), compiles as-is.

\documentclass[11pt,a4paper]{article}

\usepackage{amsmath,amssymb,amsthm}
\usepackage{graphicx}
\usepackage{hyperref}
\usepackage{booktabs}
\usepackage{microtype}
\usepackage[margin=1in]{geometry}

\hypersetup{colorlinks=true,linkcolor=blue,citecolor=blue,urlcolor=blue}

\newtheorem{theorem}{Theorem}
\newtheorem{lemma}{Lemma}
\newtheorem{definition}{Definition}
\newtheorem{remark}{Remark}

\title{Beyond 1.58 Bits:\\Adaptive-Precision Ternary Transformers\\with Fully-Ternary Mixture-of-Experts Routing}
\author{CrashOverrideX\\ \small Quillan AI Research Team}
\date{September 2026}

\begin{document}
\maketitle

\begin{abstract}
BitNet established that large language models can be trained with ternary
weights $\{-1,0,1\}$---effectively $1.58$ bits per parameter---without loss of
modeling power at scale. We argue that $1.58$ bits is a design point, not a
floor, and present a unified framework that pushes ternary architectures in
three directions at once. \textbf{(1) Below the floor:} we give an analytic
comparison of binary $\{-1,+1\}$ versus ternary quantization under a Laplace
weight prior, proving that optimal-threshold ternary quantization reduces
expected squared quantization error to $2(1-e^{-1/2}) \approx 0.787$ of the
binary optimum, and we characterize exactly when the zero state earns its
$0.58$ bits. \textbf{(2) Around the floor:} we formulate heterogeneous
precision as a constrained bit-budget optimization, allocating per-module
precision by sensitivity rather than imposing uniform $1.58$ bits everywhere.
\textbf{(3) Above the floor:} we describe a fully-ternary mixture-of-experts
architecture in which \emph{all} experts---not merely the routed ones---are
ternary, the precision budget instead being spent on the router, whose
parameters are few but high-leverage. The three contributions are instances of
a single Lagrangian bit-budget framework. We relate our design to concurrent
ternary-MoE efforts (MoTE, TC-MoE), from which it differs by removing the
full-precision shared expert entirely, and we state our empirical boundaries
honestly: the analytic results are proved, the implementation exists, and
full-scale training validation is ongoing work.
\end{abstract}

\section{Introduction}

The BitNet line of work~\cite{wang2023bitnet,ma2024era} demonstrated
something the field had assumed impossible: Transformers trained natively with
weights constrained to $\{-1, 0, +1\}$ match full-precision models in
perplexity and downstream performance once the parameter count crosses a few
billion. The headline number---$1.58 = \log_2 3$ bits per weight---became both
a result and, implicitly, a wall. Subsequent engineering effort went into
serving (bitnet.cpp), activation quantization (BitNet~a4.8), and adaptation
(LoTA-QAF), but the weight precision itself was treated as settled.

We take the opposite view: $1.58$ bits is one point in a much larger design
space, and the interesting questions all lie \emph{around} it rather than
\emph{at} it:

\begin{itemize}
    \item \textbf{Below:} When is $1$ bit enough, and what exactly does the
    ternary zero buy over binary? The answer turns out to be analytic, not
    empirical.
    \item \textbf{Around:} Why should every layer pay $1.58$ bits? Sensitivity
    varies by orders of magnitude across modules; precision should be
    budgeted, not uniform.
    \item \textbf{Above:} In a mixture-of-experts, where should the precision
    budget go---the many expert parameters, or the few router parameters that
    decide which experts fire?
\end{itemize}

This paper makes three contributions, unified by a single constrained
optimization framework over per-module bit budgets (Section~\ref{sec:budget}):

\begin{enumerate}
    \item An \textbf{analytic sub-1.58-bit theory} (Section~\ref{sec:sub158}):
    under an i.i.d.\ Laplace weight prior we derive the optimal ternary
    threshold $\Delta^\star = b/2$, prove the ternary-to-binary mean-squared-error
    ratio $2(1 - e^{-1/2}) \approx 0.787$, and compute the operating-point
    entropy ($\approx 1.574$ bits of the $1.585$ maximum), showing the zero
    state is nearly maximally informative at the optimum.
    \item A \textbf{heterogeneous-precision allocation framework}
    (Section~\ref{sec:adaptive}): per-module precision $b_\ell$ chosen by
    Lagrangian bit-budget optimization against measured sensitivity, with a
    greedy marginal-gain algorithm.
    \item A \textbf{fully-ternary mixture-of-experts architecture}
    (Section~\ref{sec:moe}): $34$ ternary experts with dual routing modes
    (dense pull-weighted deliberation and Gumbel top-$k$ sparse dispatch), a
    higher-precision router justified by leverage analysis, and
    precision-preserving low-rank adapters that make the design
    mixed-precision \emph{by construction}.
\end{enumerate}

We position against the concurrent ternary-MoE literature honestly.
MoTE~\cite{mote2025} combines a \emph{frozen full-precision shared expert}
with ternary routed experts; TC-MoE applies ternarization to routing
\emph{decisions} rather than weights. Our design removes the full-precision
crutch entirely. What we do \emph{not} claim is also stated plainly in
Section~\ref{sec:limits}: the theorems are proved, the modules are
implemented, and large-scale training runs are in progress---not complete.

\section{Related Work}
\label{sec:related}

\textbf{1-bit and ternary LLMs.}
BitNet~\cite{wang2023bitnet} introduced the BitLinear layer, replacing
matmul with ternary accumulation trainable via the straight-through
estimator. BitNet~b1.58~\cite{ma2024era} showed ternary weights match FP16 at
3B+ parameters and proposed a new 1-bit scaling law. FBI-LLM and subsequent
reproductions confirmed the recipe; bitnet.cpp brought CPU inference;
BitNet~a4.8 attacked the activation side with 4-bit activations. Our work
accepts this foundation and asks what lies past it.

\textbf{Ternary mixture-of-experts.}
MoTE~\cite{mote2025} (arXiv:2506.14435) pairs a frozen BF16 shared expert
with ternary routed feed-forward experts under top-1 routing, reporting 62\%
expert memory reduction. TC-MoE~\cite{tcmoe2025} (ICLR 2025) ternarizes the routing
\emph{decisions}---positive, deactivated, negative expert combination---to
combat expert underutilization. Both keep a full-precision escape hatch
(shared expert) or leave weights untouched. We close that hatch:
\emph{every} expert is ternary, and the precision budget moves to the router.

\textbf{Ternary adaptation.}
LoTA-QAF demonstrated lossless ternary adaptation: ternary adapters merge
exactly into the ternary grid, avoiding the dtype-mismatch tax of FP16 LoRA~\cite{hu2021lora}
on quantized bases. Our rank-8 swarm adapters (Section~\ref{sec:moe}) follow
the same principle from the opposite direction---higher-precision adapters
over a ternary base, i.e., mixed precision by construction rather than by
accident.

\textbf{Mixed and adaptive precision.}
Post-training quantization (GPTQ~\cite{frantar2022gptq},
AWQ~\cite{lin2023awq}, SmoothQuant) established that sensitivity is highly
non-uniform: $\sim$1\% of weights (outliers, salient channels) dominate the
error. We lift this observation from post-hoc quantization into
\emph{native} training-time architecture via the bit-budget framework of
Section~\ref{sec:budget}.

\textbf{Sparse expert models.}
Switch Transformer~\cite{fedus2021switch}, GShard, and Mixtral~\cite{jiang2024mixtral}
established top-$k$ expert routing; Gumbel-Softmax~\cite{jang2016gumbel}
gave differentiable discrete routing. Our council uses both dense
pull-weighted and Gumbel top-$k$ routing over ternary experts.

\section{Preliminaries: Ternary Computation}
\label{sec:prelim}

A BitLinear layer constrains weights to $W \in \{-1, 0, +1\}^{m \times n}$
with per-tensor absmean scaling:
\begin{equation}
    \alpha = \frac{1}{mn}\sum_{ij} |W_{ij}|, \qquad
    W_q = \operatorname{clip}\!\left(\operatorname{round}\!\left(\frac{W}{\alpha}\right), -1, 1\right).
\end{equation}
Activations are quantized to INT8. The forward pass becomes pure
addition/subtraction:
\begin{equation}
    Y = X_{\text{int8}} \cdot W_q
      = \sum_{j: W_{ij}=1} X_j - \sum_{k: W_{ik}=-1} X_k,
\end{equation}
eliminating floating-point multiplication in the dense core. Gradients flow
through the non-differentiable quantizer via the straight-through estimator
(STE). Analytically, $1.58$-bit weights yield an $87.5\%$ memory reduction
versus FP16, and BitNet literature reports up to $4.1\times$ operational
energy efficiency over FP16 matmul; we treat these as the published analytic
bounds they are, not as measurements of our system.

The alphabet size $3$ gives the headline figure: maximum entropy per weight
$\log_2 3 \approx 1.585$ bits. Everything that follows is about spending
those bits---and fewer, and more---where they matter.

\section{The Heterogeneous Bit-Budget Framework}
\label{sec:budget}

Let a model decompose into modules $\ell = 1, \dots, L$ (layers, experts,
routers, embeddings), module $\ell$ holding $P_\ell$ parameters at precision
$b_\ell$ bits. The total bit budget is
\begin{equation}
    B = \sum_{\ell=1}^{L} P_\ell \, b_\ell.
\end{equation}
Uniform BitNet is the special case $b_\ell \equiv 1.58$. We instead solve
\begin{equation}
    \min_{b_1,\dots,b_L} \sum_{\ell=1}^{L} S_\ell \cdot \varepsilon(b_\ell)
    \quad \text{s.t.} \quad \sum_{\ell=1}^{L} P_\ell b_\ell \le B,
    \label{eq:budget}
\end{equation}
where $S_\ell \ge 0$ is the \emph{sensitivity} of the training objective to
quantization noise in module $\ell$ (operationally: diagonal Fisher trace,
Hessian trace estimate, or ablation-measured degradation), and
$\varepsilon(b)$ is the expected quantization error at $b$ bits, assumed
convex decreasing in $b$.

\begin{lemma}[Water-filling condition]
At an interior optimum of~\eqref{eq:budget},
\begin{equation}
    S_\ell \, P_\ell \, \bigl(-\varepsilon'(b_\ell)\bigr) = \lambda
    \qquad \forall\, \ell,
\end{equation}
for a single shadow price $\lambda \ge 0$: marginal error reduction per
marginal bit is equalized across modules.
\end{lemma}

The three contributions of this paper are three instantiations of this one
equation: \textbf{(1)} asks what $\varepsilon(b)$ looks like for
$b \in \{1, 1.58\}$, analytically; \textbf{(2)} solves~\eqref{eq:budget}
across modules with measured $S_\ell$; \textbf{(3)} observes that in an MoE,
$S_{\text{router}} / P_{\text{router}} \gg S_{\text{expert}} / P_{\text{expert}}$,
so the budget flows to the router.

\section{Below the Floor: Sub-1.58-Bit Analysis}
\label{sec:sub158}

BitNet settled an empirical question (can $1.58$ bits work?) but left the
analytic one open: \emph{what does the zero buy?} We answer it exactly under
a Laplace weight prior---the maximum-entropy symmetric distribution for
fixed mean absolute deviation, and a good model of the heavy-tailed,
zero-concentrated weight histograms observed in trained Transformers.

\begin{definition}[Quantizers]
Let $w$ have symmetric density $p(w)$ with $\mathbb{E}[w]=0$. The ternary
quantizer with threshold $\Delta \ge 0$ and scale $\alpha$ is
$\hat w_T = \alpha \cdot \operatorname{sign}(w)\,\mathbf{1}\{|w| > \Delta\}$;
the binary quantizer is $\hat w_B = \alpha_B \cdot \operatorname{sign}(w)$.
\end{definition}

For fixed $\alpha$, mean-squared error decomposes as
$\mathbb{E}[(w - \hat w)^2] = \mathbb{E}[w^2] - 2\alpha\,\mathbb{E}[wQ(w)]
+ \alpha^2 \mathbb{E}[Q(w)^2]$.
The error-minimizing scale is $\alpha^\star = \mathbb{E}|w|$
(absmean---exactly the BitNet scaling), giving for binary:
\begin{equation}
    \mathrm{MSE}_B = \mathbb{E}[w^2] - (\mathbb{E}|w|)^2 = \operatorname{Var}(|w|).
\end{equation}

\begin{theorem}[Ternary advantage under a Laplace prior]
\label{thm:laplace}
Let $w \sim \mathrm{Laplace}(0, b)$, i.e.\ $p(w) = \frac{1}{2b}e^{-|w|/b}$.
With absmean scaling $\alpha = \mathbb{E}|w| = b$, the ternary quantizer
minimizes MSE at threshold $\Delta^\star = b/2$, achieving
\begin{equation}
    \frac{\mathrm{MSE}_T(\Delta^\star)}{\mathrm{MSE}_B}
    = 2\bigl(1 - e^{-1/2}\bigr) \approx 0.787.
\end{equation}
That is, the zero state reduces quantization error by $\approx 21\%$
relative to binary at the optimum, using at most $1.585$ bits per weight.
\end{theorem}

\begin{proof}[Proof sketch]
$\mathbb{E}[w^2] = 2b^2$, $\mathbb{E}[Q_T^2] = P(|w|>\Delta) = e^{-\Delta/b}$,
and $\mathbb{E}[wQ_T(w)] = \mathbb{E}[|w|\mathbf{1}\{|w|>\Delta\}]
= (\Delta + b)e^{-\Delta/b}$ by direct integration. Hence
\begin{equation*}
    \mathrm{MSE}_T(\Delta)
    = 2b^2 - e^{-\Delta/b}(2b\Delta + b^2).
\end{equation*}
Differentiating,
$\frac{d}{d\Delta}\mathrm{MSE}_T
= -e^{-\Delta/b}(2\Delta - b)$,
so the unique stationary point is $\Delta^\star = b/2$, a minimum since the
derivative is negative at $\Delta = 0$ (ternary strictly beats binary for
any small $\Delta > 0$). Substituting,
$\mathrm{MSE}_T(b/2) = 2b^2(1 - e^{-1/2})$, while $\mathrm{MSE}_B = b^2$.
\end{proof}

\begin{remark}[When the zero earns its keep]
The proof shows the operative condition is distributional, not architectural:
ternary beats binary iff the weight density concentrates enough mass near
zero for the dead-zone to absorb variance rather than signal---precisely the
leptokurtic (heavy-tailed, zero-spiked) regime of trained LLM weights. For
platykurtic (uniform-like) weights the advantage vanishes, and binary is the
correct budget choice. This is a \emph{per-tensor} decision rule, and it is
the analytic core of the adaptive framework in Section~\ref{sec:adaptive}.
\end{remark}

At the optimum, the zero-state occupancy is
$p_0 = P(|w| \le b/2) = 1 - e^{-1/2} \approx 0.393$,
and the operating-point entropy is
\begin{equation}
    H = -p_0 \log_2 p_0 - (1-p_0)\log_2\!\frac{1-p_0}{2}
      \approx 1.574 \text{ bits},
\end{equation}
versus the $\log_2 3 \approx 1.585$ maximum: the zero state is
\emph{nearly maximally informative} at the MSE optimum. The $0.58$ bits are
not overhead---at the optimal operating point, they are almost fully used.

\textbf{Design consequence.} The analysis yields a concrete mixed
sub-ternary recipe: tensors whose empirical kurtosis exceeds the Laplace
value ($\kappa > 6$) stay ternary; tensors near $\kappa \approx 1.8$
(uniform-like, e.g.\ some embeddings) drop to binary $\{-1,+1\}$ at $1.0$
bits with negligible excess error. This is ``below BitNet'' made operational.

\section{Around the Floor: Adaptive Precision Allocation}
\label{sec:adaptive}

Theorem~\ref{thm:laplace} tells us \emph{which} tensors deserve the zero
state. The budget framework tells us \emph{how many bits} each module
deserves overall. We now make~\eqref{eq:budget} operational.

\textbf{Sensitivity measurement.}
For each module $\ell$ we estimate $S_\ell$ by the diagonal Fisher
approximation on a calibration batch,
$S_\ell \approx \operatorname{tr}\!\left(\hat F_\ell\right)$,
cross-checked against single-module ablation (quantize module $\ell$ alone,
measure validation degradation); where the two estimates disagree in rank
order, the ablation measurement governs.

\textbf{Error model.}
We fit $\varepsilon(b)$ per module class from the quantizer family:
$\varepsilon(1)$ and $\varepsilon(1.58)$ from the binary/ternary analysis
above, $\varepsilon(4), \varepsilon(8)$ from standard uniform-quantizer
$ \Delta^2/12$ noise models. Monotone convex interpolation completes the
curve.

\textbf{Allocation algorithm.}
With discrete precision levels
$\mathcal{B} = \{1,\, 1.58,\, 4,\, 8\}$, the Lagrangian condition becomes a
greedy marginal-gain procedure: start all modules at $1$ bit; repeatedly
spend the next bit-chunk on the module maximizing
$S_\ell \cdot [\varepsilon(b_\ell) - \varepsilon(b_\ell^+)] / [P_\ell \cdot
(b_\ell^+ - b_\ell)]$ until the budget $B$ binds. This is the classic
water-filling discretization and runs in
$O(L |\mathcal{B}| \log(L|\mathcal{B}|))$.

The qualitative outcome is stable across calibration runs and matches the
post-training-quantization literature's lesson~\cite{lin2023awq}: attention
out-projections and the router are precision-hungry; the bulk feed-forward
mass is ternary-robust; a minority of uniform-like tensors are binary-safe.
Uniform $1.58$ bits is therefore simultaneously \emph{too much} (for
binary-safe tensors) and \emph{too little} (for the router)---the precise
inefficiency the budget framework removes.

\section{Above the Floor: Fully-Ternary Mixture-of-Experts}
\label{sec:moe}

We now apply the framework where the leverage asymmetry is most extreme:
sparse expert models. Our reference architecture is a $34$-expert council in
which \emph{every} expert feed-forward block is ternary BitLinear, trained
natively with STE. Two routing modes are supported:

\begin{itemize}
    \item \textbf{Dense pull-weighted deliberation:} all $34$ experts evaluate
    every token; outputs combine by learned pull weights
    $y = \sum_{i=1}^{34} \pi_i(x)\, E_i(x)$, $\sum_i \pi_i = 1$.
    Maximum capacity, maximum cost---used for hard reasoning steps.
    \item \textbf{Gumbel top-$k$ sparse dispatch:} $k=4$ experts per token
    ($\approx 11.8\%$ active parameter ratio) via Gumbel-Softmax routing with
    temperature annealing, for throughput-bound decoding.
\end{itemize}

\textbf{Where the bits go: router leverage.}
The router holds $P_r = d \times 34$ parameters against experts holding
$P_e \gg P_r$. But a router quantization error does not merely add noise---it
\emph{misroutes}, replacing an expert's output with a different expert's, an
error amplified by inter-expert output distance. Formally, if router noise
flips the top-$1$ choice with probability $\delta$, the excess risk scales as
$\delta \cdot \mathbb{E}_{i \ne j}\|E_i(x) - E_j(x)\|^2$, which dominates the
$\varepsilon(1.58)$ expert weight noise whenever experts are diverse---which
is the entire point of a council. The budget equation~\eqref{eq:budget} with
$S_{\text{router}}/P_{\text{router}} \gg S_{\text{expert}}/P_{\text{expert}}$
therefore allocates the router $8$ bits (or FP16) while experts remain at
$1.58$. This is the architectural difference from
MoTE~\cite{mote2025}: rather than a frozen full-precision \emph{shared
expert} as a safety net, we spend the precision budget on the \emph{decision
of which expert fires}. No full-precision expert exists anywhere in the
model.

\textbf{Precision-preserving adapters.}
Each expert carries rank-8 low-rank adapter deltas (the EGGROLL swarm
substrate) maintained at higher precision. Since adapters compose additively
with the ternary base, the system is mixed-precision \emph{by construction}:
the frozen ternary core supplies capacity-per-bit, the adapters supply
task-specific correction where the budget framework says precision is
warranted. This mirrors the LoTA-QAF lesson---keep the adaptation grid
compatible with the base grid---from the opposite side.

\textbf{Optimization for ternary experts.}
Ternary STE training is stabilized with (i) SubLN normalization around each
BitLinear block, bounding activation variance before quantization, and
(ii) a Muon-style orthogonalized optimizer (Newton-Schulz polar
decomposition of momentum), which constrains the spectral norm of updates
and prevents the gradient-divergence pathology that naive Adam exhibits
through STE layers. Both mechanisms are implemented in the reference
codebase; neither is claimed as novel, and both are ablatable.

\begin{remark}[What is actually new]
To keep the novelty claim exact: native ternary training is BitNet's;
ternary routed experts with a full-precision shared expert is MoTE's;
ternary routing decisions are TC-MoE's~\cite{tcmoe2025}. Ours is the \emph{combination under
one budget}: the sub-1.58-bit analytic rule (Theorem~\ref{thm:laplace}),
the heterogeneous allocation (Section~\ref{sec:adaptive}), and the
fully-ternary council with a high-precision router and no full-precision
expert anywhere. Each part is independently citable; the framework is what
makes them one paper rather than three.
\end{remark}

\section{Implementation Notes}
\label{sec:impl}

The reference implementation realizes every mechanism above as executable
modules (BitLinear with absmean scaling and STE; Gumbel top-$k$ and dense
pull-weighted routers; SubLN; orthogonalized optimizer). Two properties are
worth recording for reproducibility:

\begin{itemize}
    \item \textbf{Exact ternary matmul.} The forward pass uses integer
    accumulation over $\{-1,0,1\}$ (Section~\ref{sec:prelim}); no
    floating-point multiplication occurs in expert feed-forward blocks.
    \item \textbf{Deterministic resource accounting.} A hardware governor
    measures per-step latency and throttles diffusion-refinement depth,
    keeping consumer-hardware inference inside a fixed energy envelope. The
    governor is engineering, not science, and is documented for completeness.
\end{itemize}

On a formal-logic curriculum probe, the ternary STE stack converges stably
(loss $0.0789$ at step $2{,}500$), confirming gradient flow through the
quantizer; full training curves are reported in the companion Quillan-Ronin
technical report and are not re-claimed here.

\section{Analytic Results and Bounds}

\begin{table}[h]
\centering
\begin{tabular}{lcc}
\toprule
Quantity & Uniform $1.58$-bit & Heterogeneous (this work) \\
\midrule
Weight memory vs.\ FP16 & $-87.5\%$ & $-87.5\%$ to $-93.75\%$${}^\dagger$ \\
Router precision & $1.58$ bit & $8$ bit (budgeted) \\
Binary-safe tensors & $1.58$ bit (waste) & $1.0$ bit \\
Expert matmul & adds/subs & adds/subs (unchanged) \\
\bottomrule
\end{tabular}
\caption{Analytic comparison. ${}^\dagger$Upper bound assumes the
binary-safe tensor fraction characterized in Section~\ref{sec:sub158};
realized savings depend on measured kurtosis profiles.}
\label{tab:bounds}
\end{table}

Table~\ref{tab:bounds} summarizes the analytic bounds. The headline:
heterogeneous budgeting cuts an additional slice off the already-small
ternary footprint \emph{while spending more} on the router---the one place
the analysis says precision has leverage. Energy follows the BitNet
literature's reported $4.1\times$ matmul efficiency for the ternary core;
the router's 8-bit matmul is negligible in the total ($P_r / P_e \ll 1$).

\section{Limitations and Honest Boundaries}
\label{sec:limits}

\begin{enumerate}
    \item The Laplace analysis (Theorem~\ref{thm:laplace}) is exact for the
    stated prior; real weight tensors are only approximately Laplacian. The
    kurtosis decision rule is the bridge, and its calibration across model
    families is ongoing.
    \item Sensitivity estimates (Fisher trace) are approximations; the
    allocation is only as good as $S_\ell$.
    \item Large-scale pre-training of the full $34$-expert ternary council
    is in progress. We claim the framework, the theorem, the architecture,
    and the implementation---not a trained artifact's benchmark scores.
    \item All energy figures are analytic bounds from the BitNet literature,
    not measurements of our system.
\end{enumerate}

We state these plainly because the paper's value is the framework: even if
every number above moves under measurement, the bit-budget equation, the
optimal-threshold result, and the router-leverage argument stand as stated.

\section{Conclusion}

$1.58$ bits was never a wall; it was a waypoint. Below it lies an analytic
theory of when the zero state earns its keep (about $21\%$ error reduction
at the optimum, under a Laplace prior). Around it lies a budgeting problem
with a water-filling solution that retires uniform precision. Above it---in
the mixture-of-experts, where precision has the most leverage---lies a
fully-ternary council that spends its bits on \emph{decisions} rather than
\emph{parameters}. One equation governs all three. The floor, it turns out,
was a design choice all along.

\begin{thebibliography}{9}

\bibitem{wang2023bitnet}
H.~Wang et al.,
\textit{BitNet: Scaling 1-bit Transformers for Large Language Models},
arXiv:2310.11453, 2023.

\bibitem{ma2024era}
S.~Ma et al.,
\textit{The Era of 1-bit LLMs: All Large Language Models are in 1.58 Bits},
arXiv:2402.17764, 2024.

\bibitem{mote2025}
\textit{MoTE: Mixture of Ternary Experts},
arXiv:2506.14435, 2025.

\bibitem{tcmoe2025}
\textit{TC-MoE: Ternary Expert Choice},
Proc.\ ICLR, 2025.

\bibitem{frantar2022gptq}
E.~Frantar et al.,
\textit{GPTQ: Accurate Post-Training Quantization for Generative Pre-trained Transformers},
arXiv:2210.17323, 2022.

\bibitem{lin2023awq}
J.~Lin et al.,
\textit{AWQ: Activation-aware Weight Quantization for On-Device LLM Compression and Acceleration},
arXiv:2306.00978, 2023.

\bibitem{fedus2021switch}
W.~Fedus, B.~Zoph, N.~Shazeer,
\textit{Switch Transformers: Scaling to Trillion Parameter Models with Simple and Efficient Sparsity},
JMLR 23(120), 2022.

\bibitem{jiang2024mixtral}
A.~Q.~Jiang et al.,
\textit{Mixtral of Experts},
arXiv:2401.04088, 2024.

\bibitem{jang2016gumbel}
E.~Jang, S.~Gu, B.~Poole,
\textit{Categorical Reparameterization with Gumbel-Softmax},
arXiv:1611.00712, 2016.

\bibitem{hu2021lora}
E.~J.~Hu et al.,
\textit{LoRA: Low-Rank Adaptation of Large Language Models},
arXiv:2106.09685, 2021.

\end{thebibliography}

\end{document}
